\subsection{The remaining character comparison can be realized by elliptic curves}

\begin{lemma}[Integral conjugacy of the finite norm-one action]
\label{lem:B-RD35-conjugacy}
Let $k$ be algebraically closed of characteristic three, let
$E_0/k$ be $y^2=x^3-x$, and write $G_0=E_0[3^\infty]$.
Every finite reduced-norm-one subgroup of $\operatorname{Aut}(G_0)$
is conjugate by an automorphism of $G_0$ into
$\operatorname{Aut}(E_0)$. In particular this is integral conjugacy,
not only conjugacy by a quasi-isogeny.
\end{lemma}
\begin{proof}
Use the quaternion division algebra $D$ and its valuation ring
$\mathcal O_D=\End(G_0)$ from (B-RD31.1), normalizing the division
valuation by $v_D(3)=2$. Write
\[
 \zeta:(x,y)\longmapsto(x+1,y),\qquad
 j:(x,y)\longmapsto(-x,iy),\quad i^2=-1.
\]
These automorphisms have
\[
 \zeta^3=1,\quad j^2=-1,\quad
 j\zeta j^{-1}=\zeta^{-1},\qquad
 H_0=\langle\zeta,j\rangle\quad\hbox{of order twelve}.
                                                        \tag{B-RD35.1}
\]
Here $-1$ denotes the elliptic involution, not the identity.
Every origin-preserving automorphism of this curve has the form
$(x,y)\mapsto(u^2x+r,u^3y)$, with $u^4=1$ and $r^3-r=0$:
the pole orders of $x,y$ give the usual Weierstrass change of
variables, and comparison of the $y$-terms eliminates additions
to $y$ since $2$ is invertible. Substitution then gives the two
displayed conditions. Thus $H_0=\operatorname{Aut}(E_0)$.
Put $E=\mathbf Q_3(\zeta)\subset D$. We have
\[
 v_D(\zeta-1)=1,\qquad
 (\zeta-1)j(\zeta-1)^{-1}=-\zeta j.
                                                        \tag{B-RD35.2}
\]
The first equality follows from
$(\zeta-1)(\zeta^{-1}-1)=3$. Consequently $\zeta-1$
normalizes $H_0$ and permits adjustment of a conjugator's valuation
by any integer.

By Lemma~\ref{lem:B-RD32-norm}, a finite norm-one group $H$ is
cyclic of order $1,2,3,4$, or $6$, or is the nonabelian group of
order twelve described there. We give conjugators rather than
appealing only to rational Skolem--Noether.

If $z\in H$ has order three, the element $d=z+\zeta+1$
satisfies $dz=\zeta d$. If $d=0$, then $z=\zeta^{-1}$,
and $j$ instead conjugates $z$ to $\zeta$. Otherwise $d$ is
invertible because $D$ is a division algebra. Left multiplication
by a power of $\zeta-1$, which commutes with $\zeta$, makes $d$
a valuation-zero element. It is then an element of
$\mathcal O_D^\times=\operatorname{Aut}(G_0)$.
This handles the cyclic groups of order three and six; the unique
involution is $-1$.

For a generator $t$ of a cyclic group of order four, $t^2=-1$.
Unless $t=-j$, the nonzero element $d=t+j$ satisfies $dt=jd$.
Multiplication on the left by a suitable power of $3$ makes its
valuation either zero or one. In the latter case multiply once by
$q=\zeta-\zeta^{-1}$, which has valuation one and conjugates
$j$ to $-j$, and then divide by $3$. The resulting conjugator is
a unit and sends $\langle t\rangle$ onto $\langle j\rangle$.
The case $t=-j$ already has the desired subgroup.

Finally suppose $|H|=12$. First use the integral conjugacy just
proved to put its order-three subgroup inside $\langle\zeta\rangle$.
Its order-four generator $t$ conjugates $\zeta$ to $\zeta^{-1}$,
so $t=bj$ for some $b\in E^\times$. The equality $t^2=-1$ gives
$b\bar b=1$. If $b\ne-1$, take $c=1+b$; then $c/\bar c=b$.
If $b=-1$, take $c=\zeta-\zeta^{-1}$ instead, with the same
identity. Thus $c^{-1}tc=j$, while conjugation by $c$ fixes
$\zeta$. Multiplying this rational conjugator on the left by a
power of $\zeta-1$ makes it a unit without changing the target
group $H_0$, by (B-RD35.2). This proves integral conjugacy in
every case.
\end{proof}

\begin{proposition}[B-RD35: elliptic realization with its exact Tate lattice]
\label{prop:B-RD35}
\status{Proved. Independent parent and obstruction/crystal proof, original-source,
and exact-code audits passed.}
Let $K$ be a complete discretely valued field of mixed characteristic
$(0,3)$ with algebraically closed residue field. Let $M/K$ be a
height-two $3$-divisible group that has good reduction to a
one-dimensional supersingular formal group over a finite extension.
Assume its finite special-inertia action has reduced norm one.
Then there is an elliptic curve $E/K$ and an isomorphism of
$3$-divisible groups
\[
                             E[3^\infty]\simeq M.
                                                        \tag{B-RD35.3}
\]
Thus the comparison is integral on the full Tate lattice. It is not
merely an equality of rational representations or of finite inertia
characters.

In particular, after the completed unramified base change already
used in B-RD31--B-RD34, every remaining level-$27$ pair in the
$p=3,e_K=2,h=2$ problem is a pair of elliptic curves with
isomorphic $27$-torsion. The two curves have different tame special
inertia orders, norm-one special inertia, and joint faithful image
of order $18$ or $36$.
\end{proposition}
\begin{proof}
Choose a finite Galois good-reduction extension $L/K$, with
$R=\mathcal O_L$ and group $\Gamma$. Its residue extension is
trivial. Let $\mathcal G/R$ be the good model of $M_L$.
Tate's integral full-faithfulness theorem extends the generic
semilinear descent isomorphisms, and their inverses, to
$\mathcal G$; uniqueness preserves the cocycle relations
\cite[Theorem~4 and Corollaries~1--2, p.~181]{Tate1967PDivisible}.
Reduction gives the specified finite action of $\Gamma$ on
$\mathcal G_k$.

Lazard's height classification identifies $\mathcal G_k$ with
$G_0=E_0[3^\infty]$
\cite[Theorem~IV and its proof, pp.~270--271]{Lazard1955Formal}.
By Lemma~\ref{lem:B-RD35-conjugacy}, this identification can be
chosen so that the reduced action is the $3$-divisible-group action
of a homomorphism $\Gamma\to\operatorname{Aut}(E_0)$.
The homomorphism, including its labelling by elements of $\Gamma$,
is retained; no alteration of the inertia representation is made.

For $R_n=R/\pi^nR$, the prime three is nilpotent. Serre--Tate
equivalence gives a compatible elliptic deformation $\mathcal E_n/R_n$
of $E_0$, together with an isomorphism
$\mathcal E_n[3^\infty]\simeq\mathcal G_{R_n}$ respecting the
chosen special-fibre identification
\cite[Theorem~1.2.1, printed p.~143]{Katz1981SerreTate}.
Compatibility follows from uniqueness in that equivalence.
Each identity section is an effective relative Cartier divisor,
and the line bundles $\mathcal O_{\mathcal E_n}(3\cdot0)$ are
compatible and relatively ample. Grothendieck existence therefore
algebraizes the formal system to a projective curve
$\mathcal E/R$; the identity and group operations algebraize by
formal full faithfulness. The formal flatness and smoothness
criteria give smoothness along the closed fibre. Since $R$ is
local and $\mathcal E$ is proper, the closed nonsmooth locus
cannot avoid the closed fibre unless it is empty. Thus
$\mathcal E$ is an elliptic scheme over $R$. The exact existence
and full-faithfulness passages used here are the same as in C.1,
namely \cite[Theorems~5.4.1 and~5.4.5]{EGAIII1}.
For every $r$, the compatible finite-level Hopf algebra
isomorphisms modulo $\pi^n$ pass to an isomorphism over $R$,
because the coordinate algebras of $\mathcal G[3^r]$ and
$\mathcal E[3^r]$ are finite free and $R$ is complete.
They commute with the transitions in $r$. Hence
\[
                    \mathcal E[3^\infty]\simeq\mathcal G
                                                        \tag{B-RD35.4}
\]
integrally over $R$.

For every $\gamma\in\Gamma$, the specified special automorphism
of $E_0$ and the corresponding semilinear isomorphism of
$\mathcal G$ lift, by Serre--Tate full faithfulness on each $R_n$,
to a semilinear isomorphism of $\mathcal E_n$. The lifts are
compatible in $n$, and their cocycle relations hold by the same
uniqueness. Proper formal full faithfulness algebraizes these
isomorphisms and relations over $R$.

Restrict this descent datum to the generic elliptic curve
$\mathcal E_L$. Finite Galois descent is effective: the identity
section and its canonical ample line bundle
$\mathcal O(3\cdot0)$ are preserved, giving polarized projective
descent. We obtain an elliptic curve $E/K$. The isomorphism
(B-RD35.4) respects the descent data, so its generic restriction
descends to (B-RD35.3). Only the generic elliptic curve is
descended here; the ramified good model is not claimed to descend
to an elliptic scheme over $\mathcal O_K$.

The last assertion now follows from Theorem~\ref{thm:B-RD34}.
Conversely, such elliptic pairs are a subclass of the original
$3$-divisible-group question. Thus existence of a remaining
comparison after this unramified completion is equivalent to
existence in the elliptic subclass, with the same exact torsion
isomorphism and the same finite inertia characters.
\end{proof}

\begin{remark}[An explicit bounded-dimensional parameterization]
Since $2$ is invertible, every curve in this reduction has an
integral Weierstrass equation
\[
        y^2=x^3+a_2x^2+a_4x+a_6,\qquad a_2,a_4,a_6\in\mathcal O_K,
                                                        \tag{B-RD35.5}
\]
with nonzero discriminant, after scaling. Hence the remaining
question concerns two triples of integral coefficients and an
isomorphism of their rank-$3^6$ generic group schemes of
$27$-torsion. The proposition justifies this reduction for every
allowed potentially supersingular BT pair; it does not evaluate
the resulting congruence problem or assert that such a pair exists.
\end{remark}

\begin{remark}[Exact quaternion checks]
The retained program \path{code/verify_elliptic_inertia_conjugacy.py}
works in the rational quaternion algebra
$\mathbf Q(\zeta_3)\oplus\mathbf Q(\zeta_3)j$.
It verifies the displayed normalizer and conjugator identities,
all $144$ products of the twelve automorphisms, their element-order
distribution, and the integral-normalization algorithm on eighty
exact rational conjugates. The author run passed in
$0.15318646800005808$ seconds. These finite checks illustrate the
explicit algebraic proof of all cases; they do not replace its
classification or its valuation argument.
\end{remark}
